\chapter{Espaces de Banach et théorèmes fondamentaux}
\label{ch:b3:banach}

L'analyse fonctionnelle étudie les espaces normés de dimension
infinie à travers les opérateurs et fonctionnelles qui y
vivent. Sa découverte fondatrice est que la \emph{complétude},
via le théorème de Baire, force une forte uniformité~: les
familles d'opérateurs ponctuellement bornées sont bornées en
norme (Banach--Steinhaus), les bijections \omterm{def:b3:topology:continuity}{continues} ont des
inverses \omterm{def:b3:topology:continuity}{continues} (application ouverte), et les graphes
détectent la continuité (graphe fermé). L'autre pilier,
\emph{Hahn--Banach}, n'a besoin d'aucune complétude ---
seulement du lemme de Zorn --- et garantit que les espaces
duaux sont assez riches pour voir chaque vecteur. Ce chapitre
prouve les quatre théorèmes et les teste sur les espaces de
suites classiques $\ell^p$, sur des calculs de duaux concrets,
et sur une application véritablement surprenante~: il existe
des fonctions \omterm{def:b3:topology:continuity}{continues} $2\pi$-périodiques dont la série de
Fourier \emph{diverge} en un point --- résolvant par la
négative une question que la deuxième année avait laissée
ouverte.

Tout au long, $E, F$ sont des espaces normés sur $K = \R$ ou
$\C$~; \emph{espace de Banach}\index{espace de Banach} signifie
espace normé \omterm{def:b3:complete:complete}{complet}.

\section{Opérateurs bornés~; espaces de suites}

\begin{definition}\label{def:b3:banach:operator}
$\mathcal L(E, F)$ désigne l'espace des applications linéaires
\emph{bornées} (= \omterm{def:b3:topology:continuity}{continues}, deuxième année) muni de la
\emph{norme d'opérateur}\index{norme d'opérateur} $\vertiii T
= \sup_{\norm x \leq 1}\norm{Tx}$~; elle est
sous-multiplicative~: $\vertiii{ST} \leq \vertiii
S\,\vertiii T$. Le \emph{dual}\index{espace dual} est $E' =
\mathcal L(E, K)$.
\end{definition}

\begin{proposition}\label{prop:b3:banach:llcomplete}
Si $F$ est un espace de Banach, $\mathcal L(E, F)$ l'est aussi~;
en particulier $E'$ est toujours un espace de Banach.
\end{proposition}

\begin{proof}
Soit $(T_n)$ de Cauchy pour $\vertiii\cdot$. Pour chaque $x$,
$\norm{T_nx - T_mx} \leq \vertiii{T_n - T_m}\norm x$~:
$(T_nx)$ est de Cauchy dans $F$, convergente~; appeler la
limite $Tx$. $T$ est linéaire (limites d'identités linéaires)~;
en passant à la limite dans $\norm{T_nx - T_mx} \leq
\varepsilon\norm x$ ($n, m \geq N$) on obtient $\norm{T_nx -
Tx} \leq \varepsilon\norm x$~: $T_n \to T$ en \omterm{def:b3:banach:operator}{norme
d'opérateur}, et $\vertiii T \leq \vertiii{T_N} + \varepsilon <
\infty$.
\end{proof}

\begin{definition}\label{def:b3:banach:lp}
Les espaces de suites classiques\index{espaces de suites
$\ell^p$} (sur $K$, indexés par $\N$)~:
\[
\ell^p = \Bigl\{x = (x_n) : \norm x_p = \Bigl(\sum_n
\abs{x_n}^p\Bigr)^{1/p} < \infty\Bigr\}\quad (1 \leq p <
\infty),
\]
\[
\ell^\infty = \{x : \norm x_\infty = \sup\abs{x_n} < \infty\},
\]
et $c_0 = \{x : x_n \to 0\}$ muni de $\norm\cdot_\infty$. Que
$\norm\cdot_p$ soit une norme suit de l'inégalité de Minkowski,
prouvée dans le cas discret exactement comme dans
le~\cref{ch:b3:lp} (ou en sommant l'inégalité de dimension
finie de deuxième année). Tous sont des espaces de Banach, et
$c_0$ est un sous-espace fermé de $\ell^\infty$
(\cref{exo:b3:banach:3}).
\end{definition}

\begin{omfigure}
\begin{tikzpicture}[scale=1.5]
  \draw[->] (-1.45,0) -- (1.45,0) node[right,font=\small]
    {$x_1$};
  \draw[->] (0,-1.45) -- (0,1.45) node[above,font=\small]
    {$x_2$};
  % p = 1 diamond
  \draw[omDef, thick] (1,0) -- (0,1) -- (-1,0) -- (0,-1)
    -- cycle;
  % p = 2 circle
  \draw[omProp, thick] (0,0) circle (1);
  % p = 4 superellipse |x|^4 + |y|^4 = 1
  \draw[omThm, thick, smooth]
    plot[domain=0:90, samples=46]
    ({cos(\x)/(cos(\x)^4 + sin(\x)^4)^0.25},
     {sin(\x)/(cos(\x)^4 + sin(\x)^4)^0.25});
  \draw[omThm, thick, smooth]
    plot[domain=90:180, samples=46]
    ({cos(\x)/(cos(\x)^4 + sin(\x)^4)^0.25},
     {sin(\x)/(cos(\x)^4 + sin(\x)^4)^0.25});
  \draw[omThm, thick, smooth]
    plot[domain=180:270, samples=46]
    ({cos(\x)/(cos(\x)^4 + sin(\x)^4)^0.25},
     {sin(\x)/(cos(\x)^4 + sin(\x)^4)^0.25});
  \draw[omThm, thick, smooth]
    plot[domain=270:360, samples=46]
    ({cos(\x)/(cos(\x)^4 + sin(\x)^4)^0.25},
     {sin(\x)/(cos(\x)^4 + sin(\x)^4)^0.25});
  % p = infty square
  \draw[black!60, thick] (-1,-1) rectangle (1,1);
  \node[omDef, font=\small] at (0.40,0.40) {$p{=}1$};
  \node[omProp, font=\small] at (-0.52,0.75) {$p{=}2$};
  \node[omThm, font=\small] at (0.90,-0.52) {$p{=}4$};
  \node[black!60, font=\small] at (-1.62,1.15)
    {$p{=}\infty$};
  \draw[black!60, ->] (-1.32,1.11) -- (-1.02,1.0);
\end{tikzpicture}

{\small Boules unités des $p$-normes dans le plan, emboîtées
quand $p$ croît de $1$ (losange) à $2$ (disque) et $4$
(superellipse) jusqu'à $\infty$ (carré). La convexité de
chaque boule \emph{est} l'inégalité de Minkowski~; les coins
en $p = 1$ et $p = \infty$ sont là où la stricte convexité,
l'unicité des meilleures approximations, et les cas d'égalité
de l'\cref{exo:b3:lp:12} dégénèrent d'un coup.}
\end{omfigure}

\begin{proposition}[Série de Neumann]\label{prop:b3:banach:neumann}
Soit $E$ de Banach et $T \in \mathcal L(E) = \mathcal L(E,E)$
avec $\vertiii T < 1$. Alors $I - T$ est inversible dans
$\mathcal L(E)$, avec $(I - T)^{-1} = \sum_{n\geq0}T^n$
(convergente en \omterm{def:b3:banach:operator}{norme d'opérateur}). Par conséquent l'ensemble
des opérateurs inversibles est \omterm{def:b3:topology:topology}{ouvert}, et l'inversion y est
\omterm{def:b3:topology:continuity}{continue}.\index{série de Neumann}
\end{proposition}

\begin{proof}
La série converge absolument ($\vertiii{T^n} \leq \vertiii
T^n$, géométrique) dans le Banach $\mathcal L(E)$
(\cref{prop:b3:banach:llcomplete}~;
\cref{exo:b3:complete:1}(b)). En télescopant, $(I -
T)\sum_{n \leq N}T^n = I - T^{N+1} \to I$, et de même de
l'autre côté. Pour l'ouverture~: si $S$ est inversible et
$\vertiii{H} < 1/\vertiii{S^{-1}}$, alors $S + H = S(I +
S^{-1}H)$ avec $\vertiii{S^{-1}H} < 1$~: inversible.
Continuité de l'inversion~: l'expression en série donne
$\vertiii{(S+H)^{-1} - S^{-1}} = O(\vertiii H)$ localement.
\end{proof}

\begin{example}[Une équation de Volterra, résolue par Neumann]
\label{ex:b3:banach:volterraexample}
Sur $E = \mathcal C(\intcc01)$, considérer l'équation
intégrale
\[
u(x) = 1 + \lambda\int_0^xu(t)\,\dd t,
\qquad\text{c'est-à-dire}\qquad u = \mathbf 1 + \lambda Tu,
\quad (Tu)(x) = \int_0^xu .
\]
Ici $\vertiii T \leq 1$, donc pour $\abs\lambda < 1$ la \omterm{prop:b3:banach:neumann}{série
de Neumann} s'applique directement~: $u = (I - \lambda
T)^{-1}\mathbf 1 = \sum_n\lambda^nT^n\mathbf 1$. En calculant,
$T^n\mathbf 1 = \frac{x^n}{n!}$, donc
\[
u(x) = \sum_{n\geq0}\frac{(\lambda x)^n}{n!} =
\eu^{\lambda x} ,
\]
ce que la dérivation confirme. Mieux~: $\vertiii{T^n} \leq
\frac1{n!}$ (le noyau itéré se contracte factoriellement),
donc $\sum\lambda^nT^n$ converge pour \emph{tout} $\lambda$
--- l'opérateur $I - \lambda T$ est inversible pour tout
$\lambda \in \C$, même si $\vertiii{\lambda T} \geq 1$
éventuellement~: ce qui compte est la décroissance spectrale
des puissances, pas la première norme. Les opérateurs de
Volterra ont cette décroissance factorielle intégrée
(\cref{exo:b3:complete:4}(a) l'exploitait exactement), c'est
pourquoi les problèmes de Cauchy ne souffrent jamais des
phénomènes de résonance des problèmes aux limites
(\cref{ch:b3:spectral}).
\end{example}

\section{Hahn--Banach}

\begin{theorem}[Hahn--Banach, forme analytique]
\label{thm:b3:banach:hahnbanach}
Soit $E$ un espace vectoriel \emph{réel}, $p \colon E \to \R$
\emph{sous-linéaire} ($p(x + y) \leq p(x) + p(y)$ et $p(tx) =
tp(x)$ pour $t \geq 0$), $F \subseteq E$ un sous-espace et $f
\colon F \to \R$ linéaire avec $f \leq p$ sur $F$. Alors $f$
s'étend en une $\tilde f \colon E \to \R$ linéaire avec
$\tilde f \leq p$ sur $E$.\index{théorème de Hahn--Banach}
\end{theorem}

\begin{proof}
\emph{Extension en une étape.} Soit $x_0 \notin F$~; on étend
$f$ à $F \oplus \R x_0$ en choisissant $\alpha = \tilde
f(x_0)$ correctement~: on a besoin, pour tous $y \in F$, $t >
0$,
\[
f(y) + t\alpha \leq p(y + tx_0)
\quad\text{et}\quad
f(y) - t\alpha \leq p(y - tx_0),
\]
ce qui après division par $t$ (sous-linéarité) se réduit à
\[
\sup_{v \in F}\ \bigl[f(v) - p(v - x_0)\bigr]
\;\leq\; \alpha \;\leq\;
\inf_{u \in F}\ \bigl[p(u + x_0) - f(u)\bigr].
\]
Un tel $\alpha$ existe ssi tout membre de gauche est $\leq$
tout membre de droite~: en effet $f(v) + f(u) = f(u + v) \leq
p(u + v) \leq p(u + x_0) + p(v - x_0)$, c'est-à-dire $f(v) -
p(v - x_0) \leq p(u + x_0) - f(u)$.

\emph{Zorn.} Ordonner les extensions de $f$ dominées par $p$
(couples~: sous-espace, fonctionnelle) par extension~; une
chaîne a pour majorant l'union~; un élément \omterm{def:b3:rings:primemaximal}{maximal} doit être
défini sur tout $E$, sinon l'extension en une étape contredit
la maximalité.
\end{proof}

\begin{corollary}\label{cor:b3:banach:hbnormed}
Soit $E$ un espace normé ($K = \R$ ou $\C$).
\begin{enumerate}
  \item Toute $f \in F'$ ($F$ un sous-espace) s'étend en
        $\tilde f \in E'$ avec $\norm{\tilde f}_{E'} = \norm
        f_{F'}$.
  \item Pour tout $x \neq 0$ il existe $f \in E'$ avec $\norm
        f = 1$ et $f(x) = \norm x$. En particulier $E'$ sépare
        les points de $E$, et $\norm x = \sup_{\norm f \leq
        1}\abs{f(x)}$.
  \item Pour un sous-espace fermé $F$ et $x \notin F$, il
        existe $f \in E'$ s'annulant sur $F$ avec $f(x) =
        d(x, F)$ et $\norm f \leq 1$.
\end{enumerate}
\end{corollary}

\begin{proof}
(1) Cas réel~: appliquer \cref{thm:b3:banach:hahnbanach} avec
$p(x) = \norm f_{F'}\,\norm x$ (sous-linéaire)~; l'extension
satisfait $\pm\tilde f(x) = \tilde f(\pm x) \leq p(x)$, donc
$\norm{\tilde f} \leq \norm f$, et $\geq$ est la restriction.
Cas complexe~: soit $u = \operatorname{Re}f$, une fonctionnelle
réelle avec $\abs u \leq \norm f\norm\cdot$~; noter $f(x) =
u(x) - \iu\,u(\iu x)$ (vérifier sur les parties réelle et
imaginaire~: $\operatorname{Im}f(x) =
-\operatorname{Re}f(\iu x)$). Étendre $u$ $\R$-linéairement
avec la même borne, et poser $\tilde f(x) = \tilde u(x) -
\iu\tilde u(\iu x)$~: $\C$-linéaire (vérification directe sur
la multiplication par $\iu$), prolonge $f$~; norme~: pour $x$
donné écrire $\tilde f(x) = r\eu^{\iu\theta}$, alors
$\abs{\tilde f(x)} = \tilde f(\eu^{-\iu\theta}x) = \tilde
u(\eu^{-\iu\theta}x) \leq \norm f\,\norm x$.

(2) Sur $F = Kx$ définir $f(tx) = t\norm x$~: norme $1$ sur
$F$~; étendre par (1). La formule de dualité~: $\leq$ est
claire, $\geq$ par cette $f$.

(3) Sur $F \oplus Kx$ définir $f(y + tx) = t\,d(x, F)$~; alors
pour $t \neq 0$, $\norm{y + tx} = \abs t\,\norm{x + y/t} \geq
\abs t\,d(x, F) = \abs{f(y + tx)}$~: $\norm f \leq 1$ sur le
sous-espace~; étendre par (1).
\end{proof}

\begin{remark}\label{rem:b3:banach:bidual}
Par (2), l'application canonique $J \colon E \to E''$,
$J(x)(f) = f(x)$, est une isométrie
(\cref{exo:b3:banach:10})~: tout espace normé s'assied dans son
bidual. Les espaces avec $J$ surjective sont dits
\emph{réflexifs}\index{espace réflexif}~; le problème de fin de
semaine montre que $\ell^p$ ($1 < p < \infty$) est réflexif
tandis que $\ell^1$ ne l'est pas.
\end{remark}

\section{La trilogie de Baire}

\begin{theorem}[Banach--Steinhaus, bornitude uniforme]
\label{thm:b3:banach:banachsteinhaus}
Soit $E$ un espace de \emph{Banach}, $F$ normé, et
$(T_i)_{i\in I} \subseteq \mathcal L(E, F)$ une famille avec
$\sup_i \norm{T_ix} < \infty$ pour tout $x \in E$. Alors
$\sup_i \vertiii{T_i} < \infty$.\index{théorème de
Banach--Steinhaus}
\end{theorem}

\begin{proof}
Les ensembles $F_n = \{x : \sup_i\norm{T_ix} \leq n\}$ sont
fermés (intersections d'images réciproques de boules fermées)
et recouvrent $E$. Baire (\cref{thm:b3:complete:baire}) donne
$n_0$ et une boule $B(x_0, r) \subseteq F_{n_0}$. Pour $\norm z
< r$~: $\norm{T_iz} \leq \norm{T_i(x_0 + z)} +
\norm{T_ix_0} \leq 2n_0$, donc $\vertiii{T_i} \leq 2n_0/r$ pour
tout $i$.
\end{proof}

\begin{corollary}\label{cor:b3:banach:bslimits}
Si $E$ est de Banach et $T_n \in \mathcal L(E,F)$ convergent
\emph{ponctuellement} ($T_nx \to Tx$ pour chaque $x$), alors
$\sup_n \vertiii{T_n} < \infty$, $T \in \mathcal L(E, F)$, et
$\vertiii T \leq \liminf \vertiii{T_n}$.
\end{corollary}

\begin{proof}
Les suites convergentes sont bornées~: bornitude ponctuelle~;
Banach--Steinhaus borne les normes par quelque $M$~; alors
$\norm{Tx} = \lim\norm{T_nx} \leq M\norm x$ ($T$ est linéaire
comme limite ponctuelle), et la borne plus fine en passant au
$\liminf$ dans $\norm{T_nx} \leq \vertiii{T_n}\norm x$.
\end{proof}

\begin{theorem}[Séries de Fourier divergentes]
\label{thm:b3:banach:fourierdiverge}
Il existe des fonctions \omterm{def:b3:topology:continuity}{continues} $2\pi$-périodiques $f$ dont
la série de Fourier diverge en $0$~: $\sup_N\abs{S_N(f)(0)} =
\infty$. Mieux~: de telles $f$ forment une \omterm{def:b3:topology:interior}{partie dense} de
$\mathcal C(S^1)$.\index{série de Fourier, divergence}
\end{theorem}

\begin{proof}
Travailler dans $E = (\mathcal C(S^1), \norm\cdot_\infty)$, un
espace de Banach, avec les fonctionnelles $\Lambda_N(f) =
S_N(f)(0) = \frac1{2\pi}\int_{-\pi}^{\pi}f(t)\,D_N(t)\,\dd t$
(noyau de Dirichlet, deuxième année). Chaque $\Lambda_N$ est
\omterm{def:b3:topology:continuity}{continue} avec
\[
\norm{\Lambda_N} = \frac1{2\pi}\int_{-\pi}^\pi\abs{D_N(t)}\,\dd
t \;=\; L_N .
\]
($\leq$ est clair~; $\geq$~: prendre $f$ \omterm{def:b3:topology:continuity}{continue}, $\norm
f_\infty \leq 1$, approximant $\operatorname{sign}D_N$ --- le
signe a un nombre fini de sauts~; lisser chaque saut sur un
intervalle de longueur $\varepsilon$ change l'intégrale de
$O(N\varepsilon)$.) Les \emph{constantes de Lebesgue} $L_N$
tendent vers l'infini~:
\[
L_N = \frac1{2\pi}\int_{-\pi}^{\pi}
\frac{\abs{\sin\bigl((N{+}\tfrac12)t\bigr)}}{\abs{\sin(t/2)}}
\,\dd t
\geq \frac{2}{\pi}\int_0^{\pi}
\frac{\abs{\sin\bigl((N{+}\tfrac12)t\bigr)}}{t}\,\dd t
= \frac{2}{\pi}\int_0^{(N+\frac12)\pi}\frac{\abs{\sin
u}}{u}\,\dd u
\]
(en utilisant $\abs{\sin(t/2)} \leq t/2$ sur $[0, \pi]$, puis
en substituant $u = (N + \tfrac12)t$). En découpant en
arches~:
\[
\int_{(k-1)\pi}^{k\pi}\frac{\abs{\sin u}}u\,\dd u
\geq \frac1{k\pi}\int_{(k-1)\pi}^{k\pi}\abs{\sin u}\,\dd u =
\frac{2}{k\pi},
\qquad\text{donc}\qquad
L_N \geq \frac{4}{\pi^2}\sum_{k=1}^N\frac1k
\xrightarrow[N\to\infty]{} \infty .
\]
Si toute $f$ \omterm{def:b3:topology:continuity}{continue} avait $\sup_N \abs{\Lambda_N(f)} <
\infty$, Banach--Steinhaus forcerait $\sup_N\norm{\Lambda_N} <
\infty$~: contradiction. Donc quelque $f$ --- en fait un
ensemble non \omterm{rem:b3:complete:meagre}{maigre}, \omterm{def:b3:topology:interior}{dense} de $f$ (le complémentaire de
$\bigcup_M\{f: \sup_N\abs{\Lambda_Nf}\leq M\}$, une réunion
dénombrable de fermés qui, n'ayant pas d'intérieur par ce qui
précède appliqué dans toute boule, est \omterm{rem:b3:complete:meagre}{maigre}) --- a
$\sup_N\abs{S_Nf(0)} = \infty$.
\end{proof}

\begin{theorem}[Application ouverte]\label{thm:b3:banach:openmapping}
Soient $E, F$ des espaces de Banach et $T \in \mathcal L(E,
F)$ \emph{surjective}. Alors $T$ est ouverte~: $T(B_E(0,1))
\supseteq B_F(0, c)$ pour quelque $c > 0$. Par conséquent un
opérateur borné bijectif entre espaces de Banach a un inverse
borné.\index{théorème de l'application ouverte}
\end{theorem}

\begin{proof}
Écrire $B = B_E(0,1)$. La surjectivité donne $F = \bigcup_n
\overline{T(nB)} = \bigcup_n n\,\overline{T(B)}$~; Baire
(\cref{thm:b3:complete:baire}) donne un intérieur à
$\overline{T(B)}$~: quelque $B_F(y_0, 4c) \subseteq
\overline{T(B)}$. Recentraliser en $0$~: pour $\norm y < 4c$,
$y_0 + y$ et $y_0$ sont tous deux des limites d'images $Tu_k$,
$Tv_k$ avec $u_k, v_k \in B$, donc $y = \lim T(u_k - v_k)$
avec $u_k - v_k \in 2B$~: $B_F(0, 4c) \subseteq
\overline{T(2B)}$, c'est-à-dire $B_F(0, 2c) \subseteq
\overline{T(B)}$.

\emph{Retirer l'\omterm{def:b3:topology:interior}{adhérence}} (ici entre la complétude de $E$)~:
soit $\norm y < c$. Choisir $x_1 \in \frac12 B$ avec $\norm{y
- Tx_1} < c/2$ ($B_F(0,2c) \subseteq \overline{T(B)}$ mis à
l'échelle $\frac12$)~; inductivement $x_k \in 2^{-k}B$ avec
$\norm{y - T(x_1 + \dots + x_k)} < c\,2^{-k}$. La série
$\sum x_k$ converge absolument dans le Banach $E$, vers $x \in
B$ (norme $< \sum 2^{-k} = 1$), et $Tx = y$ par continuité~:
$B_F(0, c) \subseteq T(B)$. L'ouverture de $T$ sur des
\omterm{def:b3:topology:topology}{ouverts} arbitraires suit par translation et mise à l'échelle~;
pour le corollaire, l'ouverture de $T$ signifie que $T^{-1}$
est \omterm{def:b3:topology:continuity}{continue}.
\end{proof}

\begin{corollary}[Normes équivalentes]\label{cor:b3:banach:equivnorms}
Si un espace vectoriel est \omterm{def:b3:complete:complete}{complet} pour deux normes
comparables ($\norm\cdot_a \leq C\norm\cdot_b$), les normes
sont équivalentes.
\end{corollary}

\begin{proof}
L'identité $(E, \norm\cdot_b) \to (E, \norm\cdot_a)$ est
bornée et bijective entre espaces de Banach~: son inverse est
borné.
\end{proof}

\begin{theorem}[Graphe fermé]\label{thm:b3:banach:closedgraph}
Soient $E, F$ de Banach et $T \colon E \to F$ linéaire. Si le
graphe $\Gamma = \{(x, Tx)\}$ est fermé dans $E \times F$
(c'est-à-dire $x_n \to x$ et $Tx_n \to y$ impliquent $y =
Tx$), alors $T$ est borné.\index{théorème du graphe fermé}
\end{theorem}

\begin{proof}
$E \times F$ muni de $\norm{(x,y)} = \norm x + \norm y$ est
de Banach~; $\Gamma$, un sous-espace fermé, est de Banach. La
projection $\pi_E\colon \Gamma \to E$ est bornée et bijective,
donc son inverse $x \mapsto (x, Tx)$ est borné
(\cref{thm:b3:banach:openmapping})~: $\norm{Tx} \leq
\norm{(x, Tx)} \leq C\norm x$.
\end{proof}

\begin{method}\label{met:b3:banach:trilogy}
Quand atteindre quel théorème. \emph{Hahn--Banach}~: pour
produire une fonctionnelle au comportement prescrit (normer un
vecteur, s'annuler sur un sous-espace, étendre depuis un
sous-espace) --- pas de complétude nécessaire.
\emph{Banach--Steinhaus}~: pour convertir de l'information
ponctuelle en bornes uniformes --- typiquement pour montrer
qu'une opération limite est \omterm{def:b3:topology:continuity}{continue}, ou (contraposée) pour
prouver une \emph{divergence} pour quelque élément, comme pour
les séries de Fourier. \emph{Application ouverte / graphe
fermé}~: pour obtenir la continuité gratuitement à partir de
la bijectivité \omterm{def:b3:galois:algebraic}{algébrique} ou d'une propriété de fermeture du
graphe --- usage typique~: comparer deux normes complètes, ou
prouver une continuité automatique. Les trois théorèmes de
Baire exigent la complétude \emph{de la source}~;
contre-exemples sinon (\cref{exo:b3:banach:7}).
\end{method}

\section{Espaces duaux, concrètement}

\begin{theorem}\label{thm:b3:banach:duals}
Isométriquement~: $(c_0)' \cong \ell^1$ et $(\ell^1)' \cong
\ell^\infty$, via le crochet $\langle x, y\rangle = \sum_n
x_ny_n$.\index{duaux des espaces de suites}
\end{theorem}

\begin{proof}
On prouve $(c_0)' \cong \ell^1$~; la seconde identification
est l'\cref{exo:b3:banach:5}. À $y \in \ell^1$ associer
$\Lambda_y(x) = \sum x_ny_n$ ($x \in c_0$)~: absolument
convergente, avec $\abs{\Lambda_y(x)} \leq \norm
x_\infty\norm y_1$, donc $\norm{\Lambda_y} \leq \norm y_1$.
Réciproquement soit $\Lambda \in (c_0)'$~; poser $y_n =
\Lambda(e_n)$ ($e_n$ les suites unitaires). Pour tout $N$,
tester $x^{(N)} = \sum_{n \leq N}
\operatorname{sign}(\overline{y_n})\,e_n \in c_0$ (norme
$\leq 1$~; dans le cas complexe utiliser des facteurs
unimodulaires $\bar y_n/\abs{y_n}$)~: $\Lambda(x^{(N)}) =
\sum_{n\leq N}\abs{y_n} \leq \norm\Lambda$. Donc $y \in
\ell^1$ avec $\norm y_1 \leq \norm\Lambda$. Enfin $\Lambda =
\Lambda_y$~: les deux coïncident sur les $e_n$, donc sur les
suites finies, \omterm{def:b3:topology:interior}{denses} dans $c_0$ (troncation~: $\norm{x -
\sum_{n \leq N}x_ne_n}_\infty = \sup_{n > N}\abs{x_n} \to 0$
précisément parce que $x_n \to 0$)~; des fonctionnelles
\omterm{def:b3:topology:continuity}{continues} coïncidant sur un \omterm{def:b3:topology:interior}{dense} sont égales. La
correspondance est linéaire, bijective, et isométrique
($\norm{\Lambda_y} = \norm y_1$ d'après les deux inégalités).
\end{proof}

\section{Exercices}

\begin{exercise}[$\star$]\label{exo:b3:banach:1}
Calculer les \omterm{def:b3:banach:operator}{normes d'opérateur}~:
(a) les décalages $S(x_1, x_2, \dots) = (0, x_1, x_2, \dots)$
et $S^*(x_1, x_2, \dots) = (x_2, x_3, \dots)$ sur $\ell^2$~;
(b) l'opérateur de multiplication $M_a x = (a_nx_n)$ sur
$\ell^2$, pour $a \in \ell^\infty$~;
(c) la fonctionnelle $\Lambda(f) = \int_0^{1/2}f -
\int_{1/2}^1f$ sur $\mathcal C(\intcc01)$ --- montrer
$\norm\Lambda = 1$ et que la norme n'est \emph{pas} atteinte.
\end{exercise}

\begin{exercise}[$\star$]\label{exo:b3:banach:2}
Soit $E$ de Banach, $T \in \mathcal L(E)$ inversible, et $S$
avec $\vertiii{S - T} < 1/\vertiii{T^{-1}}$. Montrer que $S$
est inversible et estimer $\vertiii{S^{-1} - T^{-1}}$.
Application~: si un système linéaire $Tx = b$ est solvable
avec $T$ inversible, une perturbation suffisamment petite de
$T$ le garde uniquement solvable, avec une borne quantitative
sur le changement de solution.
\end{exercise}

\begin{exercise}[$\star\star$]\label{exo:b3:banach:3}
(a) Prouver que $\ell^1$, $\ell^\infty$ et $c_0$ sont des
espaces de Banach, et que $c_0$ est l'\omterm{def:b3:topology:interior}{adhérence} dans
$\ell^\infty$ de l'espace des suites finies.
(b) Montrer $\ell^p \subseteq \ell^q$ avec $\norm\cdot_q \leq
\norm\cdot_p$ pour $1 \leq p \leq q \leq \infty$, et que
l'inclusion est stricte.
\end{exercise}

\begin{exercise}[$\star\star$]\label{exo:b3:banach:4}
Soit $F \subseteq E$ un sous-espace fermé et $x \notin F$. En
utilisant \cref{cor:b3:banach:hbnormed}, prouver la formule de
dualité
\[
d(x, F) = \max\bigl\{\abs{f(x)} : f \in E',\ \norm f \leq 1,\
f\restriction_F = 0\bigr\}
\]
(noter~: un \emph{maximum}). En déduire que $F =
\bigcap\{\ker f : f \in E',\ f\restriction_F = 0\}$~: les
sous-espaces fermés sont exactement les intersections de
noyaux de fonctionnelles.
\end{exercise}

\begin{exercise}[$\star\star$]\label{exo:b3:banach:5}
Prouver $(\ell^1)' \cong \ell^\infty$ isométriquement, en
suivant le schéma du~\cref{thm:b3:banach:duals} (les suites
finies sont \omterm{def:b3:topology:interior}{denses} dans $\ell^1$).
Où l'argument casse-t-il pour $(\ell^\infty)'$~
\end{exercise}

\begin{exercise}[$\star\star$]\label{exo:b3:banach:6}
Soient $E, F, G$ normés avec $E$ de Banach, et $B \colon E
\times F \to G$ bilinéaire, \omterm{def:b3:topology:continuity}{continue} en chaque variable
séparément. Montrer que $B$ est (conjointement) \omterm{def:b3:topology:continuity}{continue}~:
$\norm{B(x,y)} \leq C\norm x\norm y$. \emph{(Appliquer
Banach--Steinhaus à la famille $(B(\cdot, y))_{\norm y \leq
1}$.)}
\end{exercise}

\begin{exercise}[$\star\star$]\label{exo:b3:banach:7}
(a) Sur $E = \mathcal C(\intcc01)$, comparer
$\norm\cdot_\infty$ et $\norm\cdot_1$~: l'identité $(E,
\norm\cdot_\infty) \to (E, \norm\cdot_1)$ est bornée et
bijective mais son inverse est non borné. Quelle hypothèse du
\cref{cor:b3:banach:equivnorms} échoue~
(b) Exhiber une application linéaire discontinue d'un
sous-espace \omterm{def:b3:topology:interior}{dense} de $\ell^2$ vers $K$ (p.\,ex.\ sur les
suites finies), et expliquer pourquoi cela ne contredit pas le
théorème du graphe fermé.
\end{exercise}

\begin{exercise}[$\star\star\star$]\label{exo:b3:banach:8}
(Hellinger--Toeplitz) Soit $T \colon \ell^2 \to \ell^2$
linéaire (partout définie) et \emph{symétrique}~: $\langle
Tx, y\rangle = \langle x, Ty\rangle$ pour tous $x, y$, où
$\langle x, y \rangle = \sum x_n\bar y_n$. Montrer que $T$
est borné. \emph{(Graphe fermé~: si $x_k \to x$ et $Tx_k \to
z$, tester contre $y$ arbitraire.)} Morale~: les opérateurs
symétriques non bornés --- les hamiltoniens de la mécanique
quantique --- ne peuvent jamais être définis sur tout
l'espace.
\end{exercise}

\begin{exercise}[$\star\star\star$]\label{exo:b3:banach:9}
(Théorème de Pólya sur la quadrature) Pour chaque $n$, soit
$\Lambda_n(f) = \sum_{i=0}^{n} w_{i,n}f(x_{i,n})$ une règle de
quadrature sur $\mathcal C(\intcc01)$ ($x_{i,n} \in
\intcc01$, $w_{i,n} \in \R$). Montrer que $\Lambda_n(f) \to
\int_0^1f$ pour \emph{toute} $f$ \omterm{def:b3:topology:continuity}{continue} si et seulement
si~: (i) $\Lambda_n(P) \to \int_0^1P$ pour tout polynôme $P$,
et (ii) $\sup_n\sum_i\abs{w_{i,n}} < \infty$.
\emph{(Calculer $\norm{\Lambda_n}$~; utiliser
Banach--Steinhaus et Weierstrass.)} Vérifier que les règles à
poids \emph{positifs} exactes sur les constantes satisfont
(ii) automatiquement.
\end{exercise}

\begin{exercise}[$\star\star$]\label{exo:b3:banach:10}
Montrer que $J \colon E \to E''$, $J(x)(f) = f(x)$, est une
isométrie linéaire (utiliser
\cref{cor:b3:banach:hbnormed}(2)), et qu'elle est surjective
quand $\dim E < \infty$. Montrer aussi que si $E'$ est
\omterm{def:b3:galois:separable}{séparable} alors $E$ l'est. \emph{(Choisir $x_n$ presque
normant une suite \omterm{def:b3:topology:interior}{dense} de $E'$ et montrer que leur
sous-espace engendré fermé est $E$, via
\cref{cor:b3:banach:hbnormed}(3).)}
\end{exercise}

\begin{exercise}[$\star\star$]\label{exo:b3:banach:11}
(Espaces quotients) Soit $E$ un espace de Banach et $F
\subseteq E$ un sous-espace \emph{fermé}. Sur $E/F$ définir
\[
\norm{\bar x} \;=\; d(x, F) = \inf_{y\in F}\norm{x - y} .
\]
(a) Montrer que c'est une \emph{norme} bien définie sur $E/F$
(où entre la fermeture de $F$~?), et que la projection $\pi
\colon E \to E/F$ a $\vertiii\pi \leq 1$ et envoie la boule
unité ouverte sur la boule unité ouverte.
(b) Montrer que $E/F$ est \omterm{def:b3:complete:complete}{complet}. \emph{(Utiliser le critère
de séries de l'\cref{exo:b3:complete:1}(b)~: étant donné des
classes $\bar x_k$ avec $\sum\norm{\bar x_k} < \infty$,
relever chacune en $x_k \in E$ avec $\norm{x_k} \leq
\norm{\bar x_k} + 2^{-k}$ et sommer dans $E$.)}
(c) Calculer~: pour $E = c$ (suites convergentes) et $F =
c_0$, montrer $c/c_0 \cong K$ isométriquement via $\bar x
\mapsto \lim_nx_n$.
\end{exercise}

\begin{exercise}[$\star\star$]\label{exo:b3:banach:12}
(Projections bornées et sous-espaces complémentés) Soit $E$ un
espace de Banach et $P \colon E \to E$ linéaire avec $P^2 =
P$ (une projection \omterm{def:b3:galois:algebraic}{algébrique}), $V = \operatorname{im}P$, $W
= \ker P$.
(a) Supposer $P$ bornée. Montrer que $V$ et $W$ sont fermés et
$E = V \oplus W$ avec la décomposition $x = Px + (x - Px)$.
(b) Réciproquement, supposer $E = V \oplus W$ avec
\emph{les deux} $V, W$ fermés, et soit $P$ la projection sur
$V$ le long de $W$. Montrer que $P$ est bornée.
\emph{(Graphe fermé~: si $x_n \to x$ et $Px_n \to z$, alors
$z \in V$, $x_n - Px_n \to x - z \in W$, et l'unicité de la
décomposition identifie $z = Px$.)}
(c) En déduire l'\emph{équivalence}~: un sous-espace $V$
admet une projection bornée ssi il est fermé et a un
complément \omterm{def:b3:galois:algebraic}{algébrique} fermé --- et noter (sans preuve) que des
sous-espaces fermés sans cette propriété existent ($c_0$ dans
$\ell^\infty$ est l'exemple classique)~: les espaces de
Hilbert, où $V^\perp$ marche toujours
(\cref{ch:b3:hilbert}), sont l'exception, non la règle.
\end{exercise}

\section{Problème~: la dualité des espaces
\texorpdfstring{$\ell^p$}{lp}}

\begin{problem}[{Problème de fin de semaine --- $(\ell^p)' =
\ell^q$, réflexivité, et l'étrangeté de $\ell^\infty$}]
\label{pb:b3:banach:1}
Fixer $1 < p < \infty$ et soit $q$ l'exposant conjugué,
$\frac1p + \frac1q = 1$. Le crochet tout au long est $\langle
x, y\rangle = \sum_n x_ny_n$.

\medskip
\noindent\textbf{Partie I --- Hölder et Minkowski pour les
suites.}
\begin{enumerate}
  \item (Inégalité de Young) Pour $a, b \geq 0$ montrer $ab
        \leq \frac{a^p}p + \frac{b^q}q$, en utilisant la
        concavité de $\log$ ou en étudiant $t \mapsto
        \frac{t^p}p + \frac1q - t$.
  \item (Hölder) En déduire~: $\abs{\langle x, y\rangle} \leq
        \norm x_p\norm y_q$ pour $x \in \ell^p$, $y \in
        \ell^q$~; identifier le cas d'égalité.
  \item (Minkowski) En déduire l'inégalité triangulaire pour
        $\norm\cdot_p$. \emph{(Écrire $\abs{x_n + y_n}^p \leq
        \abs{x_n}\,\abs{x_n{+}y_n}^{p-1} +
        \abs{y_n}\,\abs{x_n{+}y_n}^{p-1}$ et appliquer Hölder
        à chaque terme.)}
  \item Prouver que $\ell^p$ est \omterm{def:b3:complete:complete}{complet} et que les suites
        finies y sont \omterm{def:b3:topology:interior}{denses}.
\end{enumerate}

\medskip
\noindent\textbf{Partie II --- La dualité $(\ell^p)' =
\ell^q$.}
\begin{enumerate}[resume]
  \item Pour $y \in \ell^q$, montrer que $\Lambda_y(x) =
        \langle x, y\rangle$ définit $\Lambda_y \in (\ell^p)'$
        avec $\norm{\Lambda_y} \leq \norm y_q$, et, en testant
        sur $x_n = \abs{y_n}^{q-1}\operatorname{sign}(y_n)$
        (convenablement tronqué et normalisé), que
        $\norm{\Lambda_y} = \norm y_q$.
  \item Réciproquement, étant donné $\Lambda \in (\ell^p)'$,
        poser $y_n = \Lambda(e_n)$~; montrer $y \in \ell^q$
        avec $\norm y_q \leq \norm\Lambda$ \emph{(tester sur
        des troncatures comme à la question~5 et laisser
        croître la longueur de troncature)}, et conclure
        $\Lambda = \Lambda_y$~: l'application $y \mapsto
        \Lambda_y$ est un isomorphisme isométrique $\ell^q
        \to (\ell^p)'$.
  \item En déduire que $\ell^p$ est \emph{réflexif} pour $1 <
        p < \infty$~: en composant les deux dualités, tout
        élément de $(\ell^p)''$ vient de $\ell^p$~; vérifier
        soigneusement que le composé est le $J$ canonique.
\end{enumerate}

\medskip
\noindent\textbf{Partie III --- $\ell^1$ et $\ell^\infty$ sont
des animaux différents.}
\begin{enumerate}[resume]
  \item Montrer que $\ell^p$ ($1 \leq p < \infty$) et $c_0$
        sont \omterm{def:b3:galois:separable}{séparables}, mais $\ell^\infty$ ne l'est pas.
        \emph{(Les suites indicatrices non dénombrables de
        parties de $\N$ sont deux à deux à distance $1$.)}
  \item Déduire de l'\cref{exo:b3:banach:10} que $(\ell^1)'
        \cong \ell^\infty$ mais $(\ell^\infty)' \not\cong
        \ell^1$~: $\ell^1$ n'est \emph{pas} réflexif.
        \emph{(Si $(\ell^\infty)'$ était $\ell^1$, il serait
        \omterm{def:b3:galois:separable}{séparable}, forçant $\ell^\infty$ \omterm{def:b3:galois:separable}{séparable}.)}
  \item (Une limite de Banach, explicitement) Sur
        $\ell^\infty_\R$, soit $p(x) = \limsup_n \frac{x_1 +
        \dots + x_n}{n}$. Montrer que $p$ est sous-linéaire,
        et que sur le sous-espace $c$ des suites convergentes,
        $\mathrm{LIM}(x) = \lim x$ satisfait $\mathrm{LIM}
        \leq p$. Étendre par Hahn--Banach en $\mathrm{LIM}
        \colon \ell^\infty_\R \to \R$ et montrer~:
        $\mathrm{LIM}$ est positive ($x \geq 0 \Rightarrow
        \mathrm{LIM}(x) \geq 0$), invariante par décalage
        ($\mathrm{LIM}(x_2, x_3, \dots) = \mathrm{LIM}(x)$),
        prolonge la limite, et satisfait $\liminf x \leq
        \mathrm{LIM}(x)\leq \limsup x$.
  \item Montrer qu'une telle $\mathrm{LIM}$, vue dans
        $(\ell^\infty)'$, n'est \emph{pas} de la forme
        $\Lambda_y$ pour aucun $y \in \ell^1$~; conclure à
        nouveau $(\ell^\infty)' \neq \ell^1$.
        \emph{(Évaluer sur les suites unitaires $e_n$, puis
        sur la suite constante $1$.)}
  \item Évaluer $\mathrm{LIM}$ sur $(0,1,0,1,\dots)$, et
        montrer qu'aucune extension \emph{multiplicative}
        invariante par décalage de la limite ne peut exister
        \emph{(considérer $x = (0,1,0,1,\dots)$ et $x\cdot
        Sx$ où $S$ est le décalage)}.
\end{enumerate}

\medskip
\noindent\textbf{Partie IV --- Épilogue~: pourquoi la
réflexivité compte.}
\begin{enumerate}[resume]
  \item En utilisant \cref{cor:b3:banach:bslimits} et la
        question~6, montrer que toute suite bornée de
        $\ell^p$ ($1 < p < \infty$) a une sous-suite
        $(x^{(k)})$ qui converge \emph{faiblement}~:
        $\Lambda(x^{(k)})$ converge pour toute $\Lambda \in
        (\ell^p)'$. \emph{(Extraction diagonale sur les
        dénombrablement coordonnées~; identifier la limite
        faible dans $\ell^p$ en utilisant la bornitude
        uniforme des normes et Hölder.)} Montrer par exemple
        ($e_n$ dans $\ell^1$, contre des éléments bien choisis
        de $\ell^\infty$) que cela échoue dans $\ell^1$~: la
        compacité faible est un privilège des \omterm{rem:b3:banach:bidual}{espaces
        réflexifs}.
\end{enumerate}

\medskip
\noindent\textbf{Partie V --- La \omterm{def:b3:topology:topology}{topologie} faible au travail,
et la surprise de Schur.} Écrire $x^{(k)} \rightharpoonup x$
dans un espace normé $E$ (\emph{convergence faible}) quand
$\Lambda(x^{(k)}) \to \Lambda(x)$ pour toute $\Lambda \in
E'$.
\begin{enumerate}[resume]
  \item Compléter le recensement~: montrer $(c_0)' \cong
        \ell^1$ isométriquement, par le schéma des
        questions~5--6 (que remplacent les suites tests~?).
        Assembler la chaîne $c_0 \to \ell^1 \to \ell^\infty
        \to \dots$ de duaux successifs et marquer où la
        réflexivité échoue.
  \item Montrer que toute suite faiblement convergente d'un
        espace de Banach est bornée~: voir les $x^{(k)}$ à
        travers l'immersion canonique $J$ comme
        fonctionnelles sur $E'$ et appliquer
        Banach--Steinhaus (\cref{thm:b3:banach:banachsteinhaus})
        --- sur quel espace de Banach, et pourquoi la
        complétude y est disponible~
  \item Montrer que dans $\ell^p$, $1 < p < \infty$~:
        $x^{(k)} \rightharpoonup x$ ssi
        $\sup_k\norm{x^{(k)}}_p < \infty$ et $x^{(k)}_n \to
        x_n$ pour chaque coordonnée $n$ \emph{(un sens utilise
        Banach--Steinhaus via l'immersion canonique~; pour
        l'autre, approximer $y \in \ell^q$ par des suites
        finies)}. En déduire $e_k \rightharpoonup 0$ dans
        $\ell^2$ tandis que $\norm{e_k}_2 = 1$~: les limites
        faibles peuvent perdre de la masse.
  \item Montrer que la norme est faiblement
        \omterm{def:b3:topology:continuity}{semi-continue} inférieurement~: $x^{(k)}
        \rightharpoonup x$ implique $\norm x \leq
        \liminf_k\,\norm{x^{(k)}}$ \emph{(choisir une
        fonctionnelle normante pour $x$,
        \cref{cor:b3:banach:hbnormed})}.
  \item (Radon--Riesz dans $\ell^2$) Montrer que dans
        $\ell^2$, la convergence faible jointe à la
        convergence des normes implique la convergence en
        norme \emph{(développer $\norm{x^{(k)} - x}_2^2$)}.
        Donner un contre-exemple au même énoncé sans
        l'hypothèse de norme.
  \item (Schur, étape 1) Soit $x^{(k)} \rightharpoonup 0$
        dans $\ell^1$ et supposer, par contradiction,
        $\norm{x^{(k)}}_1 \geq \delta > 0$ le long d'une
        sous-suite. Montrer d'abord que $x^{(k)}_n \to 0$
        pour chaque $n$ (quelles fonctionnelles~?), puis
        construire récursivement des indices $k_1 < k_2 <
        \cdots$ et des entiers $0 = N_0 < N_1 < N_2 <
        \cdots$ tels que la masse de $x^{(k_j)}$ se
        concentre sur le bloc $B_j =
        \intoc{N_{j-1}}{N_j}$~:
        \[
        \sum_{n \in B_j}\bigl|x^{(k_j)}_n\bigr| \geq
        \norm{x^{(k_j)}}_1 - \frac\delta{10} .
        \]
  \item (Schur, étape 2) Définir $y \in \ell^\infty$ par $y_n
        = \operatorname{sign}\bigl(x^{(k_j)}_n\bigr)$ pour $n
        \in B_j$. Montrer $\langle x^{(k_j)}, y\rangle \geq
        \norm{x^{(k_j)}}_1 - \frac{2\delta}{10} \geq
        \frac{8\delta}{10}$ et dériver une contradiction avec
        $x^{(k)} \rightharpoonup 0$. Conclure le
        \emph{théorème de Schur}~: dans $\ell^1$, les suites
        faiblement convergentes convergent en norme.
  \item En déduire que $(e_k)$ n'a pas de sous-suite
        faiblement convergente dans $\ell^1$ (sa seule limite
        candidate est $0$, coordonnée par coordonnée --- puis
        appliquer Schur), retrouvant l'échec de compacité
        faible de la question~13~; et résoudre le paradoxe
        apparent~: dans $\ell^1$ les convergences faible et en
        norme des \emph{suites} coïncident, pourtant les
        \emph{\omterm{def:b3:topology:topology}{topologies}} faible et en norme diffèrent et les
        ensembles bornés échouent encore à être
        séquentiellement faiblement compacts --- pas de
        contradiction, seulement l'échec de la réflexivité.
  \item (Tableau de synthèse) Pour $E \in \{c_0,\ \ell^1,\
        \ell^p\ (1{<}p{<}\infty),\ \ell^\infty\}$, tabuler~:
        le \omterm{def:b3:banach:operator}{dual}~; la séparabilité~; la réflexivité~; si les
        suites bornées admettent des sous-suites faiblement
        convergentes~; et une propriété signature de chaque
        espace, justifiée en une ligne depuis ce problème.
\end{enumerate}

\medskip
\noindent\textbf{Partie VI --- Compléments~: points les plus
proches, convergence moyennée, la valeur d'une limite de
Banach.}
\begin{enumerate}[resume]
  \item (Points les plus proches~: un dividende de la
        réflexivité) Soit $F$ un sous-espace fermé de
        $\ell^p$ ($1 < p < \infty$) et $x \in \ell^p$.
        Montrer que $d = \operatorname{dist}(x, F)$ est
        \emph{atteint}~: extraire d'une suite minimisante une
        sous-suite faiblement convergente (question~13),
        garder la limite faible dans $F$ en construisant, via
        Hahn--Banach, une fonctionnelle s'annulant sur $F$
        mais pas en un point hors de $F$, et conclure avec la
        question~17. Puis montrer que le privilège n'est pas
        universel~: dans $c_0$, pour $\Lambda(x) =
        \sum_n2^{-n}x_n$, prouver que $\norm\Lambda = 1$ n'est
        pas atteint sur la boule unité, établir la formule de
        distance $\operatorname{dist}(x, \ker\Lambda) =
        \abs{\Lambda(x)}$, et en déduire qu'aucun $x \notin
        \ker\Lambda$ n'a de point le plus proche dans
        l'hyperplan fermé $\ker\Lambda$.
  \item (Banach--Saks dans $\ell^2$) Soit $x^{(k)}
        \rightharpoonup 0$ dans $\ell^2_{\R}$ avec
        $\norm{x^{(k)}}_2 \leq C$. Construire une sous-suite
        $(y_j)$ avec $\abs{\langle y_i, y_j\rangle} \leq
        \frac1j$ pour tous $i < j$, et en déduire
        \[
        \Bigl\lVert\frac{y_1 + \dots + y_m}m\Bigr\rVert_2^2
        \leq \frac{C^2 + 2}m \longrightarrow 0 :
        \]
        après extraction, les moyennes de Cesàro convergent en
        \emph{norme}. Vérifier sur $(e_k)$, dont les moyennes
        ont pour norme $\frac1{\sqrt m}$~: la convergence
        faible, inutile pour la suite elle-même
        (question~16), devient convergence en norme pour les
        moyennes.
  \item (La valeur d'une limite de Banach) Soit $L$ une limite
        de Banach quelconque (question~10) et $A_mx =
        \frac1m(x + Sx + \dots + S^{m-1}x)$. Montrer $L(A_mx)
        = L(x)$ et $\liminf A_mx \leq L(x) \leq \limsup
        A_mx$~; en déduire que \emph{toutes} les limites de
        Banach coïncident sur les suites périodiques, avec
        pour valeur la moyenne sur une période --- $\frac13$
        sur $(1, 0, 0, 1, 0, 0, \dots)$, cohérent avec la
        $\frac12$ de la question~12. Puis montrer que
        l'accord échoue en général~: pour la suite par blocs
        $x$ égale à $1$ sur $\intoc{3^{j-1}}{3^j}$ pour $j$
        pair et $0$ ailleurs, montrer que les moyennes de
        Cesàro oscillent entre $\leq \frac13$ et $\geq
        \frac23$, et construire deux limites de Banach
        $L_\pm$ avec $L_-(x) \leq \frac13 < \frac23 \leq
        L_+(x)$ \emph{(étendre depuis $c \oplus \R x$ avec
        les valeurs extrêmes admissibles $\pm$~: vérifier que
        $\Lambda(y + tx) = \lim y + t\,p(x)$ est dominée par
        le sous-linéaire $p$ de la question~10)}.
\end{enumerate}
\end{problem}
